\begin{table}[!htbp]
\centering
\caption{Temuan empiris \textcite{DaSilva2025} yang menjadi titik acuan
pengukuran \emph{improvement} penelitian ini.}
\label{tab:analisis-keterbatasan-baseline}
\begin{tabularx}{\textwidth}{|X|X|X|}
\hline
\textbf{Dimensi} & \textbf{Temuan} & \textbf{Implikasi untuk Penelitian Ini} \\
\hline
Kemiripan tekstual keseluruhan &
Sekitar 85\% pertanyaan memiliki \emph{cosine similarity} $> 0{,}5$
terhadap jawaban diterima SO &
15\% pertanyaan—yang terkait domain Framework \& Library—adalah target
utama \emph{improvement} \\
\hline
Domain paling lemah &
\emph{Frameworks and Libraries}: LLM gagal karena \emph{knowledge
cutoff} dan absennya \emph{community conventions} &
GraphRAG Framework harus secara eksplisit menjawab kegagalan pada domain ini \\
\hline
Keterlacakan sumber &
LLM menghasilkan 0\% referensi ke sumber SO &
Setiap jawaban GraphRAG harus menyertakan kutipan ke SO thread \\
\hline
Dampak pada komunitas &
Penurunan signifikan aktivitas posting SO pasca-rilis ChatGPT, namun
pengguna kembali ke SO ketika LLM gagal &
Mengkonfirmasi bahwa pengetahuan komunitas SO belum dapat digantikan LLM \\
\hline
\end{tabularx}
\end{table}